我们 Rust 标准库之旅的下一站是迭代器 API。我们已经了解 C++ 中的迭代器，可以用它来\emph{迭代}，也就是逐个遍历容器中的元素，通常还会对元素执行某种操作。虽然 Rust 的迭代器与 C++ 的迭代器同名，但实现方式却大不相同。Rust 的迭代器实际上更接近 C++20 引入的\emph{范围}（range）。
在本章中，我们将了解 Rust 迭代器库的设计、构成迭代器 API 的结构体、函数与 trait，以及它们与遍历容器元素的传统方式相比有何不同。
本章将涵盖以下主要主题：
\begin{itemize}
  \item 定义与创建迭代器
  \item 消费与组合迭代器
  \item 考察迭代器的性能
\end{itemize}
\section{技术要求}
本章示例使用 Rust 1.\allowbreak{}82.\allowbreak{}0，但也可以在任何更新的版本上运行。
如果你在安装 Rust 时需要帮助，请参阅\emph{附录}。
\section{定义与创建迭代器}
\label{idx:idx_da13be50}
\label{idx:idx_03837fbb}
从本质上说，迭代器的存在是为了更轻松地对容器中的元素执行一个操作或一系列操作。这个简单的目标背后隐藏着大量复杂性。在本节中，我们将考察 Rust 中的迭代器究竟是什么，了解迭代器的特性，并学习创建迭代器的各种方式。
\subsection{什么是迭代器？}
\label{idx:idx_a2eacbc0}
在 Rust 中，迭代器就是任何实现了 \texttt{std::iter} 模块中 \texttt{Iterator} trait 的结构体。该结构体用于保存生成下一次迭代值所需的状态。例如，如果我们要遍历 \texttt{Vec} 或切片的元素，该结构体可能会存储一个指向切片的引用，以及下一个待访问元素的索引。
而 \texttt{Iterator} trait 则只暴露一个必需的方法，仅凭它就驱动了所有迭代：\texttt{next()} 方法。Rust 中的迭代围绕着这个函数的简洁性展开，它返回一个 option，其中包含容器中下一个待访问元素的值。如果容器中已没有更多元素，返回的 option 就持有 \texttt{None} 变体。
一个简短的例子可以帮助说明这一点：
\begin{codebox}[rust]
let v = vec![0, 1, 2];
let mut it = v.iter();
println!("{:?}", it.next());
println!("{:?}", it.next());
println!("{:?}", it.next());
println!("{:?}", it.next());
\end{codebox}
这里，我们构造了一个 \texttt{Vec}，然后调用 \texttt{iter()} 返回一个遍历该向量的迭代器。我们几乎从不去想 \texttt{it} 的类型，但它是专门用于遍历切片的自定义结构体：
\texttt{std::\allowbreak{}slice::\allowbreak{}Iter<\allowbreak{}'a,\allowbreak{} T>\allowbreak{}}。回想一下，\texttt{Vec} 在需要时会被强制转换为切片，而在这种
情况下，\texttt{iter()} 方法实际上存在于切片类型上。\texttt{Iter<'a, T>} 类型中的 \texttt{'a} 生命周期意味着该结构体会持有切片的一部分，因此其存活时间不能超过切片，而 \texttt{T} 当然就是切片元素的泛型类型。
接着，我们连续四次调用 \texttt{Iterator} trait 的 \texttt{next()} 方法。这段代码的输出如下：
\begin{codeboxdark}
Some(0)
Some(1)
Some(2)
None
\end{codeboxdark}
到目前为止，Rust 的迭代器看起来也许并不特别吸引人，但事情远不止于此。在实际程序中，我们不会靠一遍遍调用 \texttt{next()} 来完成任务，但重要的是要理解，这是迭代器的基本构件。明白了这一点，下面我们来看看 Rust 迭代器的四个重要特性。
\subsection{Rust 迭代器的特性}
Rust 的迭代器是惰性的、无界的、可组合的、可扩展的。下面我们逐一考察这些特性。
前面调用 \texttt{iter()} 时发生了什么？我们创建了一个新结构体，并把它绑定到可变的 \texttt{it} 变量上。但直到我们开始通过调用 \texttt{next()} 从迭代器中取值，才真正开始干活。这就是我们所说的\emph{惰性求值}。Rust 的迭代器自身什么都不做，需要由别的东西反复调用 \texttt{next()} 方法才能工作。
\label{idx:idx_32478681}
由于迭代器是惰性求值的，它们也可以是\emph{无界的}。虽然我们绝不会指望有一个无限大的容器，但并非每个迭代器都必须指向某个容器。我们已经看到，\texttt{next()} 返回 \texttt{None} 就会结束迭代，但没有任何东西阻止我们做一个永远不返回 \texttt{None} 的迭代器，用以表示一个无限序列。
事实上，我们可以轻松地做一个只用来递增计数的迭代器：
\begin{codebox}[rust]
struct Counter {
    next_value: usize,
}
impl Iterator for Counter {
    type Item = usize;
    fn next(&mut self) -> Option<usize> {
        let value = self.next_value;
        self.next_value += 1;
        return Some(value);
    }
}
\end{codebox}
我们不会做一个永远不停递增的程序，但通过下面的代码可以看到这个迭代器如预期般工作：
\begin{codebox}[rust]
let mut counter = Counter { next_value: 0 };
println!("{:?}", counter.next());
println!("{:?}", counter.next());
println!("{:?}", counter.next());
\end{codebox}
与我们前面的例子类似，它会打印出如下内容：
\begin{codeboxdark}
Some(0)
Some(1)
Some(2)
\end{codeboxdark}
但这一次，由于迭代器是无界的，我们可以按需要任意多次地继续调用 \texttt{next()} 函数，它永远不会通过返回 \texttt{None} 来结束迭代。
迭代器还是\emph{可组合的}，这意味着我们可以创建一条迭代器链，让它们依次从彼此那里取值。为此，我们需要一个结构体，它既持有一个或多个输入迭代器，又实现了 \texttt{Iterator} trait。这种结构体在 Rust 标准库中称为\emph{适配器}。在本章中我们会看到许多组合的例子。
\label{idx:idx_22ef75b1}
最后，迭代器是\emph{可扩展的}。由于迭代器 API 高度依赖 \texttt{Iterator} trait，我们可以扩展这个 trait 来为它增添新功能，正如我们在上一章用 \texttt{byteorder} crate 所做的那样。一个名为 \texttt{itertools} 的第三方 crate 正是这么做的，它为任何实现了 \texttt{Iterator} 的类型提供了大量额外方法。虽然 \texttt{itertools} 提供的功能随着时间推移已被陆续加入标准库本身，但如果你遇到棘手的迭代问题，它仍值得一看。你需要的迭代器或适配器也许正好就在里面。
现在我们已经理解了这四个特性，下面开始更深入地研究迭代器。我们先来考察创建迭代器的几种不同方式。
\subsection{创建迭代器}
创建迭代器有几种方式。我们已经看到，其中之一是调用
\texttt{iter()} 来遍历向量或切片。在整个标准库中，把迭代器的
\label{idx:idx_a5ee54f7}
创建隐藏在这样的便捷函数背后是非常常见的做法。虽然标准库中的许多容器都提供了 \texttt{iter()} 函数，但这只是一种约定；并没有某个特定的 trait 提供这个函数。
使用 \texttt{iter()} 时，约定俗成地，我们创建的是一个在遍历容器时把每个值作为不可变引用返回的迭代器。许多容器还提供 \texttt{iter\_mut()} 方法，它创建的迭代器功能类似，但返回的是可变引用。我们可以用它来修改容器内部的值。下面是一个给可变向量的每个元素加 1 的例子：
\begin{codebox}[rust]
let mut v = vec![0, 1, 2];
for value in v.iter_mut() {
    *value += 1;
}
\end{codebox}
另一种创建迭代器的方式是使用 \texttt{into\_\allowbreak{}iter(\allowbreak{})\allowbreak{}} 函数，它由
\texttt{IntoIterator} trait 提供。当在一个拥有所有权的对象（例如 \texttt{Vec} 或其他容器）上调用时，这个
方法会把容器的所有权转移给迭代器本身。此时，我们从迭代器中得到的是拥有所有权的值而非借用的值，并且由于 Rust 的借用规则，原来的容器将不再可访问。
让我们看看 \texttt{IntoIterator} 的定义，了解一下其中的原因：
\begin{codebox}[rust]
pub trait IntoIterator {
    type Item;
    type IntoIter: Iterator<Item = Self::Item>;
    fn into_iter(self) -> Self::IntoIter;
}
\end{codebox}
\label{idx:idx_07a0f181}
如我们所见，\texttt{into\_\allowbreak{}iter(\allowbreak{})\allowbreak{}} 函数接收 \texttt{self} 作为参数，而不是更常见的 \texttt{\&self} 引用参数。这意味着对象的所有权会在函数调用中转移。这是在不克隆的情况下获得拥有所有权值的迭代器的唯一方式。再看一个带类型标注的例子，就能看出 \texttt{iter()} 与 \texttt{into\_\allowbreak{}iter(\allowbreak{})\allowbreak{}} 的区别：
\begin{codebox}[rust]
let v = vec![1, 2, 3];
for x in v.iter() {
    let _:&usize=x;
    println!("{}", x);
}
for x in v.into_iter() {
    let _: usize = x;
    println!("{}", x);
}
\end{codebox}
这里，\texttt{let \_\allowbreak{}: \&\allowbreak{}usize =\allowbreak{} x} 是一个小技巧，用来标注 \texttt{x} 的类型，因为 \texttt{for} 循环中不允许直接写标注。如果标注与实际类型不匹配，编译器就会报错。
需要注意一点：如果在借用的值上调用 \texttt{into\_\allowbreak{}iter(\allowbreak{})\allowbreak{}}，它的迭代器仍会返回借用的值。这当然是因为在这种情况下迭代器无法取得值的所有权。但 \texttt{into\_\allowbreak{}iter(\allowbreak{})\allowbreak{}} 最常见的显式用法，是创建一个消费容器、产生拥有所有权的值（而非借用值）的迭代器。
我们要探讨的最后一种创建迭代器的方式是直接生成值。我们之前用 \texttt{Counter} 例子见过这种形式，但标准库中还有几个这种形式的实用迭代器：\texttt{empty()} 和 \texttt{once()} 分别产生零个和一个值，而 \texttt{repeat()} 会无限地产生同一个值。
标准库中独立迭代器里可能最有用的是 \texttt{zip()}，它接收两个迭代器作为参数，把它们的值组合成元组传递给后续阶段，一旦其中任一迭代器没有更多值可产生就停止：
\begin{codebox}[rust]
let v1 = vec![1, 2, 3];
let v2 = vec!["one", "two", "three"];
let mut it = std::iter::zip(v1.into_iter(), v2.into_iter());
println!("{:?}", it.next());
println!("{:?}", it.next());
println!("{:?}", it.next());
\end{codebox}
\label{idx:idx_d85536ab}
这个例子会打印出如下内容：
\begin{codeboxdark}
Some((1, "one"))
Some((2, "two"))
Some((3, "three"))
\end{codeboxdark}
我们已经看过几种创建迭代器的不同方式，但一直调用 \texttt{next()} 仍然显得相当乏味。因此，在下一节中，我们来看看从迭代器中取值的不同方式。
\section{消费与组合迭代器}
\label{idx:idx_1d6e552e}
\label{idx:idx_9060ecfb}
如我们所述，消费迭代器最简单的方式就是反复调用 \texttt{next()} 方法获取每个元素，直到它返回 \texttt{None}。虽然 Rust 有一些语法辅助手段，例如 \texttt{while let} 循环，但用这种方式操作迭代器仍然很麻烦。这肯定不是唯一的办法吧？下面我们来看看语言和标准库为从迭代器消费值提供了哪些选择。
\subsection{遍历迭代器}
\label{idx:idx_535d380f}
消费迭代器最常见、最方便的方式，竟然是我们在本书中见过多次的 \texttt{for} 循环。一个也许会让人惊讶的发现是：\texttt{for} 循环存在的唯一目的就是消费迭代器。既然我们已经在使用迭代器，现在就可以完整解释这个语言构造的工作原理了。
\texttt{for} 循环纯粹是语法糖，也就是说，它隐藏了一个更复杂的实现，那个实现若反复手写会很麻烦。假设我们做如下操作：
\begin{codebox}[text]
for x in v {
    println!("{}", x);
}
\end{codebox}
它实际上会展开成类似这样的代码：
\begin{codebox}[rust]
let mut it = IntoIterator::into_iter(v);
while let Some(x) = it.next() {
    println!("{}", x);
}
\end{codebox}
实际的展开要更复杂，但这个简化版本能让我们对这个语言构造有一个大致理解。这个构造的职责是把 \texttt{in} 表达式转成迭代器，然后把它整个消费掉。\texttt{for} 循环高度依赖
我们之前遇到的 \texttt{IntoIterator} trait。
在 Rust 中，如果我们写 \texttt{for x in v}，得到的 \texttt{x} 是一系列拥有所有权的值；但如果写 \texttt{for x}
\texttt{in \&v}，得到的 \texttt{x} 则是一系列借用的值或引用。这里面有一些微妙之处，
当我们研究 \texttt{for} 循环展开后的版本时，就能开始理解它们。
现在我们知道，\texttt{for x in v} 实质上等同于 \texttt{for x in}
\texttt{IntoIterator::\allowbreak{}into\_\allowbreak{}iter(\allowbreak{}v)\allowbreak{}}，而 \texttt{for x in \&\allowbreak{}v} 等同于 \texttt{for x in IntoIterator::\allowbreak{}into\_\allowbreak{}iter(\allowbreak{}\&\allowbreak{}v)\allowbreak{}}。我们还可以注意到，\texttt{for x in \&\allowbreak{}mut v} 等同于 \texttt{for x in}
\texttt{IntoIterator::\allowbreak{}into\_\allowbreak{}iter(\allowbreak{}\&\allowbreak{}mut v)\allowbreak{}}。
\label{idx:idx_3100564f}
要让这生效，这些 \texttt{IntoIterator} 实现（\texttt{T}、\texttt{\&T} 和 \texttt{\&mut T}）都需要在我们的结构体上定义。通常，\texttt{IntoIterator<\allowbreak{}\&\allowbreak{}T>\allowbreak{}} 和 \texttt{IntoIterator<\allowbreak{}\&\allowbreak{}mutT>\allowbreak{}} 分别通过调用 \texttt{iter()} 和 \texttt{iter\_mut()} 来实现，但这只是一种约定。有些容器干脆不定义 \texttt{iter\_mut()} 和 \texttt{IntoIterator<\allowbreak{}\&\allowbreak{}mutT>\allowbreak{}}，通常是因为这样做可能会损害正确性。
关于 \texttt{for} 循环还有一点很重要：\texttt{IntoIterator} 为所有实现了 \texttt{Iterator} 的类型提供了一个一揽子实现（blanket implementation），它只是原样返回该 \texttt{Iterator}，不做任何改动。每当我们把已经是迭代器的对象作为 \texttt{for} 循环的 \texttt{in} 表达式传入时，就会走这个实现。例如，当我们写 \texttt{for x in v.\allowbreak{}iter(\allowbreak{})\allowbreak{}} 时，它大致会展开成 \texttt{for x in}
\texttt{IntoIterator::\allowbreak{}into\_\allowbreak{}iter(\allowbreak{}v.\allowbreak{}iter(\allowbreak{})\allowbreak{})\allowbreak{}}。我们用 \texttt{iter()} 创建的迭代器会
原封不动地透传过去，这要归功于一揽子实现。
综合来看，下面两个循环最终是等价的：
\begin{codebox}[text]
for x in v {}
for x in v.into_iter() {}
\end{codebox}
下面这两个循环也是如此：
\begin{codebox}[text]
for x in &v {}
for x in v.iter() {}
\end{codebox}
以及这两个循环：
\begin{codebox}[rust]
for x in &mut v {}
for x in v.iter_mut() {}
\end{codebox}
\label{idx:idx_3131f3fa}
\texttt{for} 展开的另一个推论是：标准库中为大量的类型实现了 \texttt{IntoIterator}。这包括有界范围和无界范围、大多数容器类型以及若干其他类型，从而允许写出这样的代码：
\begin{codebox}[text]
for i in 0..10 {
    println!("{}", i);
}
\end{codebox}
我们现在知道，这段代码会展开成更接近这样：
\begin{codebox}[rust]
let mut it = IntoIterator::into_iter(0..10);
while let Some(i) = it.next() {
    println!("{}", i);
}
\end{codebox}
\subsection{迭代器 trait 上的函数}
\label{idx:idx_ab83be86}
当然，\texttt{for} 循环并不是消费迭代器的唯一方式。迭代器 trait 上还有几个也会消费迭代器的函数。其中许多还有有用的附带用途。例如，迭代器 trait 上的 \texttt{sum()} 函数会消费一个由整数或浮点数值构成的迭代器，并返回这些值的和。类似地，\texttt{max()} 和
\texttt{min()} 函数会在消费迭代器后返回最大或最小的元素。这些
函数依赖于我们遍历的元素可相加或可比较，并以相应的 trait 为条件。
迭代器 trait 上用于消费迭代器的最常见、最有用的函数之一是
\texttt{collect()}。\texttt{collect()} 函数会消费迭代器，并用其值创建一个新的
容器。我们要创建的容器是泛型类型，这意味着我们常常需要用类型标注或涡轮鱼运算符（turbofish）\texttt{::<} 来指定收集到的容器类型
\texttt{>}。
用我们的 \texttt{Counter} 迭代器，可以演示如何把迭代器的值收集到一个新向量中：
\begin{codebox}[rust]
let counter = Counter { next_value: 0 };
let vector = counter.take(5).collect::<Vec<usize>>();
println!("{:?}", vector);
\end{codebox}
\texttt{take()} 函数会产生一个适配器，它在返回 \texttt{None} 之前最多从迭代器中产出指定数量的元素。这段代码运行之后，我们的向量包含的值是
\texttt{[\allowbreak{}0,\allowbreak{} 1,\allowbreak{} 2,\allowbreak{} 3,\allowbreak{} 4]\allowbreak{}}。
与其他情形一样，如果我们用别的方式为编译器指定了类型，就不需要使用涡轮鱼运算符（\texttt{::< >}）。其中一种方式就是使用标注。调整一下前面的例子，我们可以像这样省略涡轮鱼：
\begin{codebox}[rust]
let vector: Vec<_> = counter.take(5).collect();
\end{codebox}
我们不需要指定向量元素的类型。只要告诉编译器我们要收集到 \texttt{Vec} 中，就足以让它推断出正确的选择。使用标注还是涡轮鱼，多半只是风格偏好问题。
\label{idx:idx_7a783e27}
有一种特殊情况：当迭代器产生的值是 \texttt{Result} 类型时。为了给我们一个可操作的编造例子，我会做一个装满 \texttt{Result} 值的 \texttt{Vec} 供我们遍历：
\begin{codebox}[rust]
let v: Vec<Result<usize, String>> = vec![Ok(1), Ok(2), Ok(3), Ok(4)];
let result = v.into_iter().collect::<Result<Vec<_>, _>>()?;
println!("{:?}", result);
\end{codebox}
这个例子会打印出一个包含 \texttt{[\allowbreak{}1,\allowbreak{} 2,\allowbreak{} 3,\allowbreak{} 4]\allowbreak{}} 的向量。这里发生了什么？当我们拥有一个由 \texttt{Result} 值构成的迭代器时，\texttt{collect()} 函数会依次检查每一个。如果迭代器中的所有值都是 \texttt{Ok}，那么每个 \texttt{Ok} 结果所持有的值会被收集到一起，并包在一个 \texttt{Ok} 中。如果其中有任何结果是 \texttt{Err}，整个过程就会停止，并返回它遇到的第一个
\texttt{Err}，作为收集到的值。
虽然收集到 \texttt{Result} 值中是种非常强大的模式，但工程师有时需要一些练习才能完全理解它。我建议你在 Rust Playground 上搭建自己的例子并不断调整，直到你彻底明白发生了什么。
既然我们已经理解如何创建和消费迭代器，下面来看看迭代器最强大的特性之一——把它们链式连接或组合起来，以高效地执行复杂操作。
\subsection{用适配器组合迭代器}
\label{idx:idx_ed06f720}
如前所述，适配器是一种特殊结构体，它既实现了迭代器 trait，又持有一个或多个作为输入来源的迭代器。标准库提供的适配器通常是迭代器 trait 的一部分，因此每个实现了该 trait 的类型都可以使用它们。
适配器的一个非常简单的例子是 \texttt{cloned()} 函数。它产生一个迭代器，该迭代器取输入迭代器的值，克隆它们，然后传给链中的下一个迭代器。当我们从引用中收集时，可能会想用它：
\begin{codebox}[rust]
let vector = v.iter().cloned().collect::<Vec<_>>();
\end{codebox}
\label{idx:idx_e628fee2}
这段代码会逐个克隆迭代器的值，然后传给 \texttt{collect()} 函数以创建新向量。更简单的做法是直接调用 \texttt{v.clone()}，但这个例子希望能说明适配器如何用来把迭代器链在一起。让我们更仔细地看看 \texttt{cloned()}，理解它在做什么。回想一下，这是实现于迭代器 trait 上的函数，意味着它可以以 \texttt{self} 身份访问调用它的迭代器。它会创建一个名为 \texttt{Cloned<I>} 的结构体，其中 \texttt{I} 是迭代器的类型。由于我们要消费该迭代器，\texttt{Cloned} 结构体只是取得迭代器的所有权；无需持有生命周期。
\texttt{Cloned} 上 \texttt{Iterator} 实现的 \texttt{next()} 函数只是利用了
\texttt{Option} 类型上的 \texttt{cloned()} 方法：
\begin{codebox}[rust]
fn next(&mut self) -> Option<T> {
    self.it.next().cloned()
}
\end{codebox}
\texttt{Option::\allowbreak{}cloned(\allowbreak{})\allowbreak{}} 方法的原理是：如果 \texttt{self} 是 \texttt{Some}，就返回一个内部含克隆值的 \texttt{Some}；如果该 option 是 \texttt{None}，就返回 \texttt{None}。只要某个适配器能用 \texttt{Option} 的方法实现，这就是一个很有用的捷径，因为迭代器的 \texttt{next()} 方法本身总是返回一个 \texttt{Option}。
真正让 Rust 迭代器 API 大放异彩的，是迭代器 trait 上提供的大量适配器。我们将介绍几个常见的。首先看 \texttt{map()} 适配器，它可能是整个库中最常用的适配器之一。\emph{map} 这个词在函数式语言中相当常见。就我们的目的而言，可以把它理解为：对迭代器返回的每个值应用一个闭包，并作为适配器把结果传给链中的下一个迭代器。
下面是一个把向量中每个数字的值翻倍并收集到新向量的例子：
\begin{codebox}[rust]
let v = vec![0, 1, 2, 3, 4];
let result = v.into_iter().map(|val| val * 2).collect::<Vec<_>>();
println!("{:?}", result);
\end{codebox}
该程序的打印输出如下：
\begin{codeboxdark}
[0, 2, 4, 6, 8]
\end{codeboxdark}
另一个常见适配器是 \texttt{filter()}，它用闭包测试每个元素。闭包返回 \texttt{true} 的元素会被传递下去，而返回 \texttt{false} 的元素会被跳过。为了只保留偶数，我们可以用 \texttt{filter()} 这样做：
\begin{codebox}[rust]
let v = vec![0, 1, 2, 3, 4];
let result = v.into_iter().filter(|val| val % 2 == 0).collect::<Vec<_>>();
println!("{:?}", result);
\end{codebox}
\label{idx:idx_025c4ede}
这个示例输出如下：
\begin{codeboxdark}
[0, 2, 4]
\end{codeboxdark}
当然，我们还可以把适配器串联起来，用一系列小步骤表达即便很复杂的计算。下面是一个同时使用 \texttt{filter()} 和 \texttt{map()} 适配器的示例：
\begin{codebox}[rust]
let v = vec![0, 1, 2, 3, 4];
let result = v
    .into_iter()
    .filter(|val| val % 2 == 0)
    .map(|val| val * 2)
    .collect::<Vec<_>>();
println!("{:?}", result);
\end{codebox}
这一次，在滤掉奇数并把剩下的值翻倍之后，我们得到如下输出：
\begin{codeboxdark}
[0, 4, 8]
\end{codeboxdark}
这里我们同时使用了 \texttt{filter()} 和 \texttt{map()}，但这种组合太常见了，因此 Rust 提供了 \texttt{filter\_\allowbreak{}map(\allowbreak{})\allowbreak{}} 适配器来完成同样的工作。\texttt{filter\_\allowbreak{}map(\allowbreak{})\allowbreak{}} 适配器返回一个 Option，其中 \texttt{Some} 值会被保留，而 \texttt{None} 值则被跳过。用 \texttt{filter\_\allowbreak{}map(\allowbreak{})\allowbreak{}} 重写的前一个示例如下：
\begin{codebox}[rust]
let v = vec![0, 1, 2, 3, 4];
let result = v
    .into_iter()
    .filter_map(|val| if val % 2 == 0 { Some(val * 2) } else { None })
    .collect::<Vec<_>>();
println!("{:?}", result);
\end{codebox}
\label{idx:idx_71df5031}
它的功能和输出都与前一个示例相同，但往往是完成同一工作更简洁、更直接的方式。花些时间翻阅一下关于迭代器 trait 的 Rust 在线文档吧。其中提供了几十个可用于生成适配器的实用方法。在任何给定的程序中，我们大概只会用到其中的少数几个，但了解有哪些可用的选项，有助于我们在自己的程序中充分利用迭代器 API。
\section{考察迭代器的性能}
\label{idx:idx_34f1c64f}
作为 C++ 程序员，我们经常要处理对性能敏感的代码。面对这些可能相当复杂的迭代器表达式链，我们会很自然地关心它们的运行时性能。但由于惰性求值、泛型函数求值以及大量内联，优化构建下的性能其实相当出色。在一些测试中，它甚至被证明能胜过手写的循环。本节将考察 Rust 编译器是如何做到这一点的。
\subsection{剖析一个迭代器表达式}
\label{idx:idx_5ae9448d}
不亲眼看到实际例子，可能很难相信迭代器的性能会胜过手写循环，所以下面来看一个简单的例子。我们将深入剖析一个迭代器表达式，手工展开并内联它，以便体会编译器眼中的它会是什么样子。
\label{idx:idx_439767f1}
我们的表达式是相对简单的 \texttt{v.\allowbreak{}into\_\allowbreak{}iter(\allowbreak{})\allowbreak{}.\allowbreak{}filter(\allowbreak{}|x| x \% 2 =\allowbreak{}=\allowbreak{} 0)\allowbreak{}.\allowbreak{}map(\allowbreak{}|}
\texttt{x| x * 2)\allowbreak{}.\allowbreak{}sum(\allowbreak{})\allowbreak{}}，其中 \texttt{v} 是一个由 \texttt{u32} 构成的向量。首先要知道的是，\texttt{sum()} 是
用 \texttt{Iterator} trait 上的 \texttt{fold()} 函数实现的。\texttt{fold()} 函数会对每个元素反复调用一个函数，每次都会传下去一个累加值。如果把 \texttt{sum()} 中所用的 \texttt{fold()} 函数展开，它看起来是这样的：
\begin{codebox}[rust]
let mut accum = 0;
while let Some(x) = self.next() {
    accum = accum + x;
}
accum
\end{codebox}
这看起来有点像我们平常会写的东西，但由于 \texttt{while let} 循环这层语法糖，很难像编译器那样看待它，所以我们来处理一种对编译器更友好的表示形式，以便看清 \texttt{next()} 函数是如何被内联的：
\begin{codebox}[rust]
let mut accum = 0;
loop {
    let x = self.next();
    if x.is_none() {
        break;
    }
    accum = accum + x.unwrap();
}
accum
\end{codebox}
\label{idx:idx_127c5aad}
\label{idx:idx_e9997999}
从 \texttt{sum()} 往回看，目标表达式中的下一项是 \texttt{map(\allowbreak{}|x| x * 2)\allowbreak{}}。由于 \texttt{map()} 是一个适配器，不弄清传入的是什么就无法展开它。在这个示例中，我们用一个占位值 \texttt{prev} 来表示。\texttt{map()} 适配器的 \texttt{next()} 函数只是通过调用 \texttt{Option::map} 实现的。不过这里我们会跳过几步，把它也内联进去，得到类似这样的结果：
\begin{codebox}[rust]
let mut accum = 0;
loop {
    let prev = self.next();
    if prev.is_none() {
        break;
    }
    let x = prev.unwrap() * 2;
    accum = accum + x;
}
accum
\end{codebox}
注意，在这种情况下，\texttt{is\_none()} 检查和 \texttt{unwrap()} 被上移到了对 \texttt{prev} 结果的处理上，而不再是 \texttt{sum()} 的一部分。这是为了便于说明而采用的简化，但在真正的编译中，优化器能够看出：用 \texttt{Some()} 创建 \texttt{Option} 之后立刻把它解包，是一个可以被消除的空操作。
在表达式中继续往回看，来到 \texttt{filter()}，它的实现稍微更有意思一些，因为当传给 \texttt{filter()} 的函数返回 \texttt{false} 时，它可以跳过某些值。\texttt{filter()} 适配器最终是用 \texttt{try\_fold()}函数实现的，从概念上讲，这会在我们手工内联的代码中加入一个嵌套循环，但可以简化成下面这样：
\begin{codebox}[rust]
let mut accum = 0;
loop {
    let value = self.next();
    if value.is_none() {
        break;
    }
    if value.unwrap() % 2 == 0 {
        let x = value.unwrap() * 2;
        accum = accum + x;
    }
}
accum
\end{codebox}
同样，为了清晰起见，这里做了一些处理。更详尽的内联版本可能会包含一个外层循环和一个内层循环，但为了便于说明，这里手工做了简化。最后一步是弄清 \texttt{v.\allowbreak{}into\_\allowbreak{}iter(\allowbreak{})\allowbreak{}} 实际上做了什么，并把它加进我们展开后的循环里。
\label{idx:idx_0f25231e}
\label{idx:idx_c86482e6}
事实证明，\texttt{Vec} 上的 \texttt{into\_\allowbreak{}iter(\allowbreak{})\allowbreak{}} 产生的是一个使用裸指针和 unsafe 块的、经过高度优化的实现。它太复杂，无法粘贴到我们这个小示例里，所以我们只用 \texttt{[]} 运算符来索引值。请记住，索引通常会在运行时触发边界检查，但在这个示例中，我们用它来替代一个更优化、不执行边界检查的实现。
这样，我们最终内联后的循环看起来就像这样：
\begin{codebox}[rust]
let mut accum = 0;
let mut start_index = 0;
let end_index = v.len();
loop {
    if start_index == end_index {
        break;
    }
    let value = v[start_index];
    if value % 2 == 0 {
        let x = value * 2;
        accum = accum + x;
    }
    start_index += 1;
}
accum
\end{codebox}
现在，这看起来可能跟我们自己动手写这个循环时的写法差不多，而且似乎能编译成还算合理的机器码。为了验证，我们可以用 \textbf{Rust Playground} 或 Matt Godbolt 的 \textbf{Compiler Explorer} 之类的工具，把这个循环优化后的汇编与原始的迭代器表达式作比较。虽然存在一些细微差异，但在这个测试中，内层循环的汇编是完全相同的。
\subsection{向量化优化}
\label{idx:idx_04dbbc57}
\label{idx:idx_031e1d31}
\label{idx:idx_c1e4d90f}
\label{idx:idx_4b0cff68}
Rust 中迭代器 API 额外做得很好的一点，是保留了\emph{向量化}的机会。向量化是 LLVM 及其他编译器后端所执行的一类优化的统称。其总体思路是：如果编译器能证明循环的某些性质，就能在现代硬件上用高效的\textbf{单指令多数据}（\textbf{SIMD}）指令来实现它。这通常意味着，数据量较大时循环的执行速度会显著提升。
众所周知，向量化很难做对，因为编译器为了证明某个特定循环可以向量化而必须遵守的各种限制很难得到保证，而且往往只有查看循环生成的汇编，才能确认你没有违反任何必需的条件。得益于其函数式设计和细致的实现，Rust 的迭代器 API 构造大多数循环的方式很可能正好满足编译器的向量化规则。这就是为什么我们有时会看到迭代器表达式比传统循环更快。
\label{idx:idx_288b3c01}
不过，Rust 迭代器虽然能让编译器更有可能执行这种优化，但这并无保证。当需要极致运行时性能时，务必运行基准测试或在本地检查生成的代码。正如前面讨论过的，Matt Godbolt 的 Compiler Explorer（\href{https://rust.godbolt.org/}{\texttt{https:/\allowbreak{}/\allowbreak{}rust.\allowbreak{}godbolt.\allowbreak{}org/\allowbreak{}}}）或 Rust Playground 等在线工具，对这类验证非常有用。
\section{小结}
至此，我们已经学习了大量关于 Rust 迭代器 API 的知识。我们发现迭代器就是实现了迭代器 trait 的结构，并学会了如何自己编写迭代器，以及如何从已有的容器创建迭代器。接着，我们考察了若干消费迭代器的方式，包括把它们收集到一个新容器中。我们还考察了适配器，以及它们如何让我们组合迭代器，构造出强大的表达式。最后，我们详细地走了一遍一个迭代器表达式的实现过程，看清了编译器为什么能够高效地优化它们。
在下一章中，我们将探讨 Rust 对最流行的编程范式之一——\textbf{面向对象}——的处理方式。
\needspace{26\baselineskip}
\begin{samepage}
\subsection{获取本书的 PDF 版本及更多内容}
扫描二维码（或访问 \texttt{packtpub.\allowbreak{}com/\allowbreak{}unlock}）。按书名搜索本书，确认版本，然后按照页面上的步骤操作。
\myGraphic{0.32}{images/qr_unlock.png}{}
\myGraphic{0.20}{images/unlock_btn.png}{}
\emph{注意：请备好你的发票。直接从 Packt 购买无需发票。}
\end{samepage}
